\section{Define an eight-by-eight prediction boundary}
Implement \texttt{h263\_predict8}. It accepts a pointer to an eight-byte-aligned
64-byte request in \texttt{\$a0} and a pointer to a sixteen-byte-aligned 192-byte
output record in \texttt{\$a1}. Each request field is an unsigned native-endian
64-bit quantity. This record describes one already positioned block; it contains
no motion-vector syntax.
\begin{center}
  \tagpdfsetup{table/header-rows={1}}
  \begin{tabular}{rll}
    \textbf{Byte offset} & \textbf{Field}    & \textbf{Meaning or range}           \\
    0                    & Current pointer   & Top-left current sample             \\
    8                    & Reference pointer & Integer origin before interpolation \\
    16                   & Current stride    & 8 through 4096 bytes                \\
    24                   & Reference stride  & $8+p_x$ through 4096 bytes          \\
    32                   & $p_x$             & 0 or 1                              \\
    40                   & $p_y$             & 0 or 1                              \\
    48                   & Mode              & 0 scalar; 1 MSA                     \\
    56                   & Reserved          & Zero
  \end{tabular}
\end{center}
The caller supplies mapped current samples in eight rows and mapped reference
samples in the required halo. Their byte pointers may be unaligned and the two
read-only inputs may overlap, including with different strides. The output is
disjoint from the inputs and request. Preserve saved general registers, the stack
pointer, the request and every input byte.

Validate record alignment and non-wrapping 64-/192-byte extents before request
reads. Then validate all fields and the minimal image extents
$7s_C+8$ and $(7+p_y)s_R+8+p_x$ before reading pixels. Use full-width field and
address arithmetic. On invalid arguments return \texttt{\$v0}=1 and
\texttt{\$v1}=0, with no pixel reads or output changes. The caller's mapped-memory
preconditions do not require the function to probe arbitrary unmapped addresses.

On success, write the following complete output and return zero in
\texttt{\$v0}. Return $\sum |e|$ in \texttt{\$v1}, an unsigned value from zero
through 16\,320. This score summarizes the residual magnitudes; it does not
replace the signed residual array.
\begin{center}
  \tagpdfsetup{table/header-rows={1}}
  \begin{tabular}{rll}
    \textbf{Byte range} & \textbf{Contents} & \textbf{Layout}                                   \\
    0--63               & Predicted samples & Eight tight rows of eight unsigned bytes          \\
    64--191             & Signed residuals  & Eight tight rows of eight native-endian halfwords
  \end{tabular}
\end{center}
A residual's byte offset is $64+2(8y+x)$. Prediction row stride is eight bytes;
residual row stride is sixteen. Those output strides are independent of both
input strides. A whole-vector halfword store fits one residual row. A whole-vector
byte load does not generally fit one eight-sample input row. Gather active bytes
into halfword lanes, or use another proven access arrangement. Coalescing adjacent
tight rows is permitted when every accessed byte belongs to the admitted footprint.

Use \path{construction/h263-prediction-construction/main.asm} as the new ordinary caller.
Its unfinished entry returns development status four. Prepare a worked copy from
your current source:
\begin{verbatim}
python3 construction/h263-prediction-construction/prepare.py \
    --source my-prediction.asm --mode 1 \
    --phase-x 1 --phase-y 1 --output worked-prediction.asm
\end{verbatim}
The helper installs Cb samples from raw pictures one and zero, using current
origin $(8,16)$ and reference origin $(7,15)$ with stride 88. It preserves every
instruction byte and refuses an existing output path. The default half-sample
phases correspond to full-luma displacement $(-1,-1)$. This data preparation
isolates interpolation; the final encoder must calculate origins itself and
use its locally reconstructed reference.

The first current sample is 128. Its four reference neighbours are 128, 128,
128 and 96, so prediction is 120 and residual is eight. The completed default
caller prints \texttt{120 8 232}, where 232 is the block's residual SAD.
At \texttt{prediction\_caller\_return}, before the caller stores returned
registers, the output record is already complete and \texttt{\$v0}=0,
\texttt{\$v1}=232. The named status/SAD fields are still awaiting those copies.
